% Part: computability
% Chapter: computability-theory
% Section: total

\documentclass[../../../include/open-logic-section]{subfiles}

\begin{document}

\olfileid{cmp}{thy}{tot}
\olsection{સંપૂર્ણતા અનિર્ણેય છે}

કોઈ વસ્તુ અસંગણનીય છે તે બતાવવા $s$-$m$-$n$ પ્રમેયનો
વધુ એક ઉપયોગ જોઈએ. $\fn{Tot}$ એ સંપૂર્ણ સંગણનીય વિધેયોના
સૂચકોનો ગણ હોય, એટલે
\[
\fn{Tot} = \Setabs{x}{\text{દરેક $y$ માટે, $\cfind{x}(y)\fdefined$}}.
\]

\begin{prop}
\ollabel{prop:total}
$\fn{Tot}$ સંગણનીય નથી.
\end{prop}

\begin{proof}
$\fn{Tot}$ સંગણનીય નથી તે જોવા $K$ને તેમાં ઘટાડી
શકાય છે એવું બતાવવું પૂરતું છે. $h(x,y)$ને આમ વ્યાખ્યાયિત કરો:
\[
h(x,y) \simeq
\begin{cases}
0 & \text{જો $x \in K$} \\
\fundefined & \text{અન્યથા}
\end{cases}
\]
નોંધો કે $h(x,y)$, $y$ પર જરાય આધાર રાખતું નથી.
$h$ આંશિક સંગણનીય છે તે સમજવું મુશ્કેલ નથી: ઇનપુટ $x, y$
પર~$h$ ગણવા પહેલાં ઇનપુટ~$x$ પર વિધેય~$\cfind{x}$નું
અનુકરણ કરીએ; તે સંગણના થંભે, તો $h(x,y)$, $0$ આપે અને
થંભે. આમ $h(x,y)$ એ $\Zero(\umin{s}{T(x,x,s)})$ જ છે,
જ્યાં $\Zero$ અચળ શૂન્ય વિધેય છે.

$s$-$m$-$n$ પ્રમેયથી એવું આદિમ પુનરાવર્તી વિધેય~$k(x)$
છે કે દરેક $x$ અને~$y$ માટે
\[
\cfind{k(x)}(y) =
\begin{cases}
0 & \text{જો $x \in K$} \\
\fundefined & \text{અન્યથા}
\end{cases}
\]
તેથી $x \in K$ હોય તો $\cfind{k(x)}$ સંપૂર્ણ છે અને
અન્યથા અનિર્ધારિત છે. આમ $k$, $K$નું~$\fn{Tot}$માં
ન્યૂનીકરણ છે.
\end{proof}

\begin{digress}
હકીકતમાં $\fn{Tot}$ સંગણકીય રીતે પરિગણનીય પણ નથી---
તેની જટિલતા ``અંકગણિતીય પદાનુક્રમ''માં વધુ ઊંચે આવે છે.
પરંતુ અહીં આ વધુ મજબૂત પરિણામની ચિંતા કરીશું નહીં.
\end{digress}

\end{document}
